\subsubsection{SVM với các \tr{kernel} PDS}

Như ta đã học, chúng ta đã thấy rằng bài toán tối ưu đối ngẫu của SVM,
cũng như dạng nghiệm của nó, không phụ thuộc trực tiếp vào các
\tr{vector} dữ liệu đầu vào mà chỉ phụ thuộc vào các phép tích trong
giữa chúng. Do \tr{kernel} PDS có khả năng ngầm định nghĩa một không
gian tích trong (theo Định lý~\ref{thm:rkhs}), ta có thể mở rộng SVM
bằng cách kết hợp với bất kỳ \tr{kernel} PDS $K$ nào thông qua việc
thay thế mỗi tích trong $x \cdot x'$ bằng hàm \tr{kernel} $K(x,x')$.

Nhờ vậy, SVM có thể hoạt động trong các không gian đặc trưng có số
chiều rất lớn, thậm chí vô hạn chiều, mà không cần biểu diễn tường
minh các \tr{vector} đặc trưng. Điều này dẫn đến dạng tổng quát của
bài toán tối ưu và nghiệm SVM với \tr{kernel} PDS như sau:
\begin{equation}
\max_{\boldsymbol{\alpha}}
\quad
\sum_{i=1}^{m}\alpha_i
-
\frac{1}{2}
\sum_{i,j=1}^{m}
\alpha_i\alpha_jy_iy_jK(x_i,x_j),
\label{eq:svm-kernel}
\end{equation}

với các ràng buộc:
\begin{equation}
0 \le \alpha_i \le C,
\quad
\sum_{i=1}^{m}\alpha_i y_i = 0,
\quad
\forall i \in [m].
\end{equation}

trong đó:
\begin{itemize}
    \item $\alpha_i$ là các biến đối ngẫu;

    \item $y_i \in \{-1,+1\}$ là nhãn lớp;

    \item $C$ là tham số điều chuẩn;

    \item $K(x_i,x_j)$ là \tr{kernel} PDS.
\end{itemize}

Sau khi giải bài toán tối ưu, hàm phân lớp thu được là:
\begin{equation}
h(x)
=
\sign\left(
\sum_{i=1}^{m}\alpha_i y_i K(x_i,x)
+ b
\right).
\label{eq:kernel-hypothesis}
\end{equation}

Hệ số chệch $b$ được tính bởi
\begin{equation}
b
=
y_i
-
\sum_{j=1}^{m}\alpha_j y_j K(x_j,x_i),
\end{equation}
với mọi $x_i$ thỏa mãn $0 < \alpha_i < C$.

Ta cũng có thể viết lại bài toán tối ưu
\eqref{eq:svm-kernel}
dưới dạng \tr{vector} bằng cách sử dụng ma trận
\tr{kernel} $\mathbf{K}$ ứng với tập huấn luyện
$(x_1,\dots,x_m)$:
\begin{equation}
\max_{\boldsymbol{\alpha}}
\quad
\mathbf{1}^{\top}\boldsymbol{\alpha}
-
\frac{1}{2}
(\boldsymbol{\alpha}\circ\mathbf{y})^{\top}
\mathbf{K}
(\boldsymbol{\alpha}\circ\mathbf{y}),
\end{equation}

với các ràng buộc
\begin{equation}
\mathbf{0}
\le
\boldsymbol{\alpha}
\le
\mathbf{C},
\quad
\boldsymbol{\alpha}^{\top}\mathbf{y}
= 0.
\end{equation}

Trong đó:
\begin{itemize}
    \item $\boldsymbol{\alpha}\in\mathbb{R}^{m\times1}$ là \tr{vector} các biến đối ngẫu;

    \item $\mathbf{y}\in\mathbb{R}^{m\times1}$ là \tr{vector} nhãn lớp;

    \item $\mathbf{K}\in\mathbb{R}^{m\times m}$ là ma trận \tr{kernel} với
    $K_{ij}=K(x_i,x_j)$;

    \item $\boldsymbol{\alpha}\circ\mathbf{y}$ là tích Hadamard (entry-wise product).
\end{itemize}

Do đó, $\boldsymbol{\alpha}\circ\mathbf{y}$ là \tr{vector} cột trong
$\mathbb{R}^{m\times1}$, với phần tử thứ $i$ bằng $\alpha_i y_i$.

Nghiệm dưới dạng \tr{vector} tương tự như trong
\eqref{eq:kernel-hypothesis}, nhưng với
\begin{equation}
b
=
y_i
-
(\boldsymbol{\alpha}\circ\mathbf{y})^{\top}
\mathbf{K}\mathbf{e}_i,
\end{equation}
với mọi $x_i$ thỏa mãn $0 < \alpha_i < C$.

Phiên bản SVM sử dụng \tr{kernel} PDS này là dạng tổng quát của SVM
sẽ được sử dụng trong các phần tiếp theo. Việc mở rộng này đặc biệt
quan trọng vì nó cho phép ánh xạ phi tuyến một cách ngầm định các điểm
dữ liệu đầu vào vào một không gian có số chiều cao hơn, nơi việc phân
tách với biên lớn (large-margin separation) có thể được thực hiện hiệu quả.

Nhiều thuật toán khác trong các lĩnh vực như regression, ranking,
dimensionality reduction hay clustering cũng có thể được mở rộng bằng
\tr{kernel} PDS theo cùng một cơ chế.